% Agent 18: PDF pages 34–35 (book pages 408–411)
% Section 5.2 Accretion in Binary Systems: The General Picture
% Covers: Roche lobe, Lagrange points, Figs 5.2.1–5.2.2, viscous stresses

\subsection{Accretion in Binary Systems: The General Picture}
\label{sec:5.2}

Consider a close binary system with one component a ``normal'' star and the
other a ``compact star'' (black hole or neutron star or white dwarf).  ``Close''
means that the separation between the centers of mass is within a factor~2 or
so of the radius $R_N$ of the normal star
\begin{equation}
a \lesssim 2\,R_N. \tag{5.2.1}
\end{equation}

For such a system, the interaction between the stars, acting over astronomical
time scales, may have brought the normal star into co-rotation with its
companion (``same face'' always turned toward companion).  For this reason,
and to simplify the discussion, we assume a circular orbit and co-rotation
of the normal star.  By Kepler's laws, the angular velocity of the stars
about each other is
\begin{equation}
\boldsymbol{\Omega} = \left[(M_N + M_c)/a^3\right]^{1/2}\,\mathbf{e}_z,
\tag{5.2.2}
\end{equation}
where $\mathbf{e}_z$ is a unit vector perpendicular to the orbital plane, and
where $M_N$ and $M_c$ are the mass of the normal star and the compact star,
respectively.  For cases of interest (e.g., the binary systems associated
with Cyg~X-1 and Cen~X-3),
\begin{gather}
M_N \sim M_c \sim 1 \text{ to } 20\,M_\odot, \notag\\
a \sim 2R_N \sim 10^{11} \text{ to } 10^{12}\;\text{cm} \tag{5.2.3}\\
\text{orbital period} = 2\pi/\Omega \sim 1 \text{ to } 5\;\text{days.}
\notag
\end{gather}

Analyze the flow of gas from the normal star to its compact companion using
a coordinate system that co-rotates with the binary system (noninertial
frame!).  In the Newtonian realm (i.e., everywhere except near the surface
of the compact star), the Euler equation~(2.5.30) for the flowing gas
reduces to
\begin{equation}
\frac{d\mathbf{v}}{dt}
  = -\boldsymbol{\nabla}\Phi_{gc}
    - \boldsymbol{\Omega}\times\mathbf{v}
    - \frac{1}{\rho_0}\,\boldsymbol{\nabla}p
    + \frac{2}{\rho_0}\,\boldsymbol{\nabla}\cdot(\eta\,\sigma).
\tag{5.2.4}
\end{equation}
Here $d/dt$ is the time derivative relative to the rotating frame:
$d/dt = \partial/\partial t + \mathbf{v}\cdot\boldsymbol{\nabla}$;
$\mathbf{v}$ is the gas velocity relative to the rotating frame; $p$ and
$\rho_0$ are pressure and rest-mass density; $\Phi_{gc}$ is the
``gravitational-plus-centrifugal'' potential
\begin{equation}
\Phi_{gc} = -\frac{M_N}{|\mathbf{r}-\mathbf{r}_N|}
            -\frac{M_c}{|\mathbf{r}-\mathbf{r}_c|}
            -\tfrac{1}{2}(\Omega\times\mathbf{r})^2;
\tag{5.2.5}
\end{equation}
and we have included a viscous shear stress, with $\eta$ the coefficient of
dynamic viscosity and $\sigma$ the shear.  In the Euler equation the term
$-\boldsymbol{\nabla}\Phi_{gc}$ gives rise to gravitational and centrifugal
accelerations; and the term $-\boldsymbol{\Omega}\times\mathbf{v}$ gives
rise to Coriolis accelerations.

Notice that when pressure and viscous accelerations are \emph{not} important
(test-particle motion), the equation of motion~(5.2.4) is identical to that
for a particle with electric charge~$q$ moving in an electric field with
potential $\Phi_{gc}/q$ and a uniform magnetic field with strength
$B = \Omega/q = (\Omega/q)\,\mathbf{e}_z$.  Conservation of energy in this
case requires $\Phi_{gc} + \tfrac{1}{2}v^2 = \text{const}$.  When pressure
forces are taken into account, energy conservation [Bernoulli
equation~(2.5.38)] is modified to read
\begin{equation}
\Phi_{gc} + \tfrac{1}{2}v^2 + w
  = \text{constant along a streamline},
\tag{5.2.6}
\end{equation}
where $w$ is the specific enthalpy (enthalpy per unit mass).

When viscous stresses are taken into account, one cannot write such a simple
conservation law.

To deduce the qualitative nature of the gas flow from the ``normal'' star to
its companion, one must have a clear picture of the potential $\Phi_{gc}$.
Figure~\ref{fig:5.2.1} gives such a picture, for the
Newtonian-plus-centrifugal potential in the orbital plane of a binary star
system with a circular orbit.

Of particular interest is the equipotential curve marked ``Roche lobe.''
If the surface of the normal star is inside its Roche lobe, then the only
way it can dump gas onto its companion is by means of a ``stellar wind,''
which blows gas off the star being captured.  But if the surface of the
normal star fills its Roche lobe, it can dump gas continuously through the
``Lagrange point''~$L_1$ onto its companion.  In this case the steady-state
flow, as governed by the Euler equation~(5.2.5) and by the law of mass
conservation,
\begin{equation}
\boldsymbol{\nabla}\cdot(\rho_0\,\mathbf{v})
  = -\partial\rho_0/\partial t = 0,
\tag{5.2.7}
\end{equation}
will have the qualitative form shown in Fig.~\ref{fig:5.2.2}.

\begin{figure}[htbp]
\centering
% \includegraphics[width=\columnwidth]{fig_5_2_1}
\fbox{\parbox{0.9\columnwidth}{\centering [Figure 5.2.1 placeholder]\\
Equipotentials of the Newtonian-plus-centrifugal potential $\Phi_{gc}$ in
the orbital plane for a binary system with mass ratio
$M_N\!:\!M_c = 10\!:\!1$.
Contour values are labelled in units of $(M_N+M_c)/a$.
Lagrange points $L_1$--$L_5$ are marked; the ``Roche lobe'' is the
figure-eight-shaped equipotential through $L_1$.  The centre of mass
(CM) is indicated.}}
\caption{The equipotentials $\Phi_{gc}=\text{const}$ for the
Newtonian-plus-centrifugal potential in the orbital plane of a binary star
system with a circular orbit.  For the case shown here the stars have a
mass ratio $M_N\!:\!M_c=10\!:\!1$.  The equipotentials are labelled by
their values of $\Phi_{gc}$ measured in units of $(M_N+M_c)/a$, where $a$
is the separation of the centres of mass of the two stars.  The innermost
equipotentials are nearly spheres.  The potential $\Phi_{gc}$ has local
stationary points ($\boldsymbol{\nabla}\Phi_{gc}=0$), called ``Lagrange
points,'' at the locations marked~$L_j$.  Of particular interest is the
equipotential curve marked ``Roche lobe.''  If the surface of the normal
star is inside its Roche lobe, then the only way it can dump gas onto its
companion is by means of a ``stellar wind,'' which blows gas off the star
being captured.  But if the surface of the normal star fills its Roche
lobe, it can dump gas continuously through the ``Lagrange point''~$L_1$
onto its companion.  In this case the steady-state flow, as governed by
the Euler equation~(5.2.5) and by the law of mass conservation,
will have the qualitative form shown in Fig.~\ref{fig:5.2.2}.}
\label{fig:5.2.1}
\end{figure}

The gas falls from the surface of the normal star, over the ``lip'' of the
potential (Lagrange point~$L_1$), toward the compact companion.  As the gas
picks up speed, Coriolis forces swing it to the right in
Figure~\ref{fig:5.2.2}, and then gravitational forces swing it back leftward
into roughly circular motion about the compact companion.  Gas previously
dumped onto the companion is now in an orbiting disk.  The incoming gas
interacts viscously with the gas of the disk.  Some of the incoming gas

\begin{figure}[htbp]
\centering
% \includegraphics[width=\columnwidth]{fig_5_2_2}
\fbox{\parbox{0.9\columnwidth}{\centering [Figure 5.2.2 placeholder]\\
Qualitative picture of the gas flow from a normal star through the inner
Lagrange point $L_1$ onto a compact companion in a close binary system.
The gas stream is deflected by Coriolis forces and forms an accretion disk
around the compact object.}}
\caption{Qualitative picture of the flow of gas from a normal star,
through the inner Lagrange point~$L_1$, onto a compact companion.
Coriolis forces deflect the incoming stream, and the gas settles into an
accretion disk orbiting the compact object.}
\label{fig:5.2.2}
\end{figure}

It is instructive to figure out why, even though $L_4$ and $L_5$ are maxima
of the potential, they are stable points for test-particle orbits.
[Particles placed at $L_4$ or $L_5$, when perturbed slightly, go into
stable orbits about $L_4$ or $L_5$.  The key to this stability is the role
of Coriolis forces; see eq.~(5.2.4).]  The ``trojan asteroids'' orbit the
Lagrange points $L_4$ and $L_5$ of the Sun--Jupiter system.
